% Gabarit pandoc des corrigés publiés (brut et guidé). Contenu sans marque ni nom de site.
% Compilation : lualatex (latexmk -lualatex).
% Le Markdown reste la source unique : ne pas éditer ce fichier .tex à la main.
\documentclass[11pt,a4paper]{article}

% --- Moteur, polices (Latin Modern) et langue ---------------------------------
\usepackage{iftex}
\ifLuaTeX\else\errmessage{Ce document se compile avec lualatex}\fi
\usepackage{amsmath,amssymb,mathtools}
\usepackage{fontspec}
% mathrm=sym : \mathrm prend ses glyphes dans la police mathématique (ℓ, μ, Ω), comme KaTeX.
\usepackage[math-style=TeX,bold-style=TeX,mathrm=sym]{unicode-math}
% Repli pour les glyphes absents de Latin Modern (indices et exposants Unicode,
% flèches, filets des schémas en texte, symboles) : DejaVu, puis Latin Modern Math
% (ℓ, lettres mathématiques Unicode des légendes).
\directlua{
  luaotfload.add_fallback("repliTexte", {"DejaVuSerif:mode=harf;", "DejaVuSans:mode=harf;", "file:latinmodern-math.otf:mode=harf;"})
  luaotfload.add_fallback("repliMono", {"DejaVuSansMono:mode=harf;", "DejaVuSans:mode=harf;", "file:latinmodern-math.otf:mode=harf;"})
}
\setmainfont{Latin Modern Roman}[RawFeature={fallback=repliTexte}]
\setsansfont{Latin Modern Sans}[RawFeature={fallback=repliTexte}]
\setmonofont{Latin Modern Mono}[RawFeature={fallback=repliMono}]
\setmathfont{Latin Modern Math}
\usepackage[french]{babel}
\frenchsetup{SmallCapsFigTabCaptions=false}

% --- Mise en page --------------------------------------------------------------
\usepackage[a4paper,margin=2.2cm,headheight=14pt,footskip=1.2cm]{geometry}
\usepackage{microtype}
\setlength{\emergencystretch}{4em}
\tolerance=1500
\setlength{\parindent}{0pt}
\setlength{\parskip}{0.55\baselineskip plus 2pt minus 1pt}
\usepackage{xcolor}
% Gris nommés : babel-french rend « ! » actif dans le document.
\definecolor{gris2}{gray}{0.98}
\definecolor{gris3}{gray}{0.97}
\definecolor{gris5}{gray}{0.95}
\definecolor{gris35}{gray}{0.65}
\definecolor{gris40}{gray}{0.60}
\definecolor{gris45}{gray}{0.55}
\definecolor{gris50}{gray}{0.50}
\definecolor{gris55}{gray}{0.45}
\definecolor{gris60}{gray}{0.40}
\definecolor{gris65}{gray}{0.35}
\definecolor{gris70}{gray}{0.30}
\definecolor{gris80}{gray}{0.20}
\definecolor{gris85}{gray}{0.15}
\usepackage{graphicx}
\usepackage{float}
\usepackage{caption}
\captionsetup{font=small,labelformat=empty,justification=centering}
\usepackage{booktabs,longtable,array,multirow,calc}
\usepackage{siunitx}
\sisetup{locale=FR}
\usepackage[version=4]{mhchem}
\usepackage{listings}
\usepackage{fancyhdr}
\usepackage{lastpage}
\usepackage[most]{tcolorbox}
\usepackage{enumitem}
\setlist{nosep,leftmargin=*,topsep=2pt}
\usepackage[hidelinks,unicode]{hyperref}
\usepackage{bookmark}
\hypersetup{pdftitle={Concours ATS — Français 2026},pdfauthor={Rédigé par IA},pdfcreator={LaTeX},pdfsubject={Corrigé guidé}}
\usepackage{xurl}

% --- Pandoc : listes, tableaux, figures ----------------------------------------
\providecommand{\passthrough}[1]{#1}
\providecommand{\tightlist}{\setlength{\itemsep}{0pt}\setlength{\parskip}{0pt}}

\makeatletter
\def\maxwidth{\ifdim\Gin@nat@width>\linewidth\linewidth\else\Gin@nat@width\fi}
\def\maxheight{\ifdim\Gin@nat@height>0.55\textheight 0.55\textheight\else\Gin@nat@height\fi}
\makeatother
\setkeys{Gin}{width=\maxwidth,height=\maxheight,keepaspectratio}
\setcounter{secnumdepth}{-\maxdimen}
\providecommand{\pandocbounded}[1]{#1}
\usepackage{longtable,booktabs,array}

% --- Code (listings) -----------------------------------------------------------
\lstset{
  basicstyle=\ttfamily\small,
  columns=fullflexible,
  keepspaces=true,
  breaklines=true,
  breakatwhitespace=false,
  postbreak=\mbox{\textcolor{gris50}{$\hookrightarrow$}\space},
  frame=single,
  rulecolor=\color{gris35},
  framesep=4pt,
  xleftmargin=5pt,xrightmargin=5pt,
  backgroundcolor=\color{gris3},
  keywordstyle=\bfseries,
  commentstyle=\color{gris65},
  stringstyle=\color{gris80},
  showstringspaces=false,
  tabsize=4,
  extendedchars=true,
  aboveskip=6pt,belowskip=6pt,
}

% --- Macros des corrigés (KaTeX ↔ LaTeX) ---------------------------------------
% Les corrigés sont écrits pour KaTeX : on définit ici ce que KaTeX connaît et
% que LaTeX n'a pas, sans jamais retoucher le Markdown.
\providecommand{\R}{\mathbb{R}}
\providecommand{\N}{\mathbb{N}}
\providecommand{\Z}{\mathbb{Z}}
\providecommand{\Q}{\mathbb{Q}}
\providecommand{\C}{\mathbb{C}}
\providecommand{\dd}{\mathrm{d}}
\providecommand{\llbracket}{\lBrack}
\providecommand{\rrbracket}{\rBrack}
\providecommand{\lBrace}{\lbrace}
\providecommand{\rBrace}{\rbrace}
\providecommand{\Box}{\square}
\providecommand{\Diamond}{\lozenge}
\providecommand{\varnothing}{\emptyset}
\providecommand{\argmax}{\operatorname*{arg\,max}}
\providecommand{\argmin}{\operatorname*{arg\,min}}
\renewcommand{\leq}{\leqslant}
\renewcommand{\geq}{\geqslant}

% Équation hors texte : réduite à la largeur de la ligne si elle déborde.
\newsavebox{\boiteequation}
\newcommand{\equationajustee}[1]{%
  \sbox{\boiteequation}{$\displaystyle #1$}%
  \ifdim\wd\boiteequation>\linewidth
    \[\resizebox{\linewidth}{!}{\usebox{\boiteequation}}\]%
  \else
    \[#1\]%
  \fi}

% --- Encadrés pédagogiques (sobres, lisibles en noir et blanc) -----------------
\tcbset{
  encadre/.style={
    enhanced,breakable,sharp corners,
    colback=white,colframe=gris55,
    boxrule=0.4pt,leftrule=2.2pt,
    left=6pt,right=6pt,top=3pt,bottom=3pt,
    before skip=6pt,after skip=6pt,
    fonttitle=\bfseries\small,coltitle=black,
    attach title to upper,after title={\quad},
    parbox=false,
  }
}
\newtcolorbox{indication}[1]{encadre,colframe=gris45,colback=gris2,title={Indication~#1}}
\newtcolorbox{methode}{encadre,colframe=gris70,title={Méthode}}
\newtcolorbox{piege}{encadre,colframe=gris85,boxrule=0.8pt,leftrule=2.2pt,title={Piège}}
\newtcolorbox{jury}{encadre,colframe=gris60,colback=gris5,title={Jury}}
\newtcolorbox{encadrenote}{encadre,colframe=gris40}

% --- Titres de section ---------------------------------------------------------
\usepackage{titlesec}
\titleformat{\section}{\Large\bfseries}{}{0pt}{}
\titleformat{\subsection}{\large\bfseries}{}{0pt}{}
\titleformat{\subsubsection}{\normalsize\bfseries}{}{0pt}{}
\titlespacing*{\section}{0pt}{14pt plus 4pt minus 2pt}{6pt}
\titlespacing*{\subsection}{0pt}{10pt plus 3pt minus 2pt}{4pt}
\titlespacing*{\subsubsection}{0pt}{8pt plus 2pt minus 2pt}{2pt}

% --- En-tête et pied de page ---------------------------------------------------
\pagestyle{fancy}
\fancyhf{}
\fancyhead[L]{\small\textsc{Corrigé guidé} · ATS ATS}
\fancyhead[R]{\small Français 2026}
\fancyfoot[L]{\footnotesize Corrigé rédigé par IA, non officiel}
\fancyfoot[R]{\footnotesize \thepage\,/\,\pageref*{LastPage}}
\renewcommand{\headrulewidth}{0.4pt}
\renewcommand{\footrulewidth}{0pt}
\fancypagestyle{plain}{\fancyhf{}\fancyfoot[R]{\footnotesize \thepage\,/\,\pageref*{LastPage}}\fancyfoot[L]{\footnotesize Corrigé rédigé par IA, non officiel}\renewcommand{\headrulewidth}{0pt}}

\begin{document}
\thispagestyle{plain}
\begin{center}
  {\small\scshape Corrigé guidé}\par\smallskip
  {\LARGE\bfseries Concours ATS — Français 2026\par}\medskip
  {\large ATS · ATS · Français · 2026\par}
\end{center}
\vspace{2pt}\hrule\vspace{8pt}

\emph{Corrigé rédigé par IA, non officiel. Signalez toute erreur.}

\hypertarget{pruxe9sentation-du-sujet}{%
\section{Présentation du sujet}\label{pruxe9sentation-du-sujet}}

\textbf{L'épreuve.} Quatre heures, sans dictionnaire ni appareil électronique, deux exercices sur un même support : un \textbf{résumé} (120 mots à 10 \% près, soit \textbf{108 à 132 mots}, une barre tous les 20 mots, total indiqué à la fin) et une \textbf{dissertation} sur une phrase extraite du texte, à traiter « en vous appuyant de manière précise sur les œuvres au programme ». L'énoncé ne donne ni barème ni définition du mot.

\textbf{Le barème (rapport 2026).} Chaque exercice est noté sur 10. Traiter un seul des deux divise la note par deux. L'orthographe coûte au plus 2 points sur 20, mais les incorrections qui rendent une phrase fautive ou incompréhensible sont comptées comme des non-sens : un résumé qui en est truffé « ne peut prétendre à une note supérieure à 1/10 ».

\textbf{Le texte.} Extrait d'un article de sciences humaines de Stéphanie Chanvallon et Stéphane Héas, « L'Homme et la Nature : en quête/enquête sensible » (\emph{Natures Sciences Sociétés}, 2011), d'environ 560 mots comptés entre blancs (environ 600 si l'on compte les élisions). Le résumé impose donc une contraction d'environ un cinquième. Le texte avance en quatre mouvements :

\begin{longtable}[]{@{}
  >{\raggedright\arraybackslash}p{(\columnwidth - 4\tabcolsep) * \real{0.0895}}
  >{\raggedright\arraybackslash}p{(\columnwidth - 4\tabcolsep) * \real{0.6344}}
  >{\raggedright\arraybackslash}p{(\columnwidth - 4\tabcolsep) * \real{0.2461}}@{}}
\toprule\noalign{}
\begin{minipage}[b]{\linewidth}\raggedright
Mouvement
\end{minipage} & \begin{minipage}[b]{\linewidth}\raggedright
Contenu
\end{minipage} & \begin{minipage}[b]{\linewidth}\raggedright
Articulation
\end{minipage} \\
\midrule\noalign{}
\endhead
\bottomrule\noalign{}
\endlastfoot
§ 1 & L'écologie, enjeu national et planétaire, oblige chacun à interroger son rapport à la Nature et son impact sur elle, comme individu et comme membre de l'espèce. & constat \\
§ 2 & Pourtant, dans les sociétés technicisées, ce rapport est épisodique, indirect, quasi artificiel ; on cherche à renouer, mais le vécu individuel reste difficile à saisir. Comme la société (Lussault), l'individu « invente » ses natures, d'autant que la mobilité et le partage des savoirs permettent des expériences ailleurs. Ces expériences touchent l'individu entier ; un rapport privilégié change le regard, \textbf{surtout} reconduit à soi, parfois transforme. & « Pourtant » ; « cela vaut aussi » ; « non seulement\ldots{} mais surtout » \\
§ 3 & En Occident, le rapport à la Nature est paradoxal (amour et exploitation) ; deux visions s'opposent, l'homme maître ou l'homme espèce parmi les autres ; la première recule, à cause de ses dégâts et du brouillage de la frontière homme / animal. L'altérité demeure, mais devient un défi à penser à travers l'animalité. & opposition, puis recul \\
§ 4 & Que faire ? Non pas un « écocitoyen » (catégorie sociopolitique), mais un habitant de la Terre qui agirait par compréhension des interdépendances dont il fait partie, non par intérêt ou bonne conscience. La Nature n'est pas uniquement ersatz, objet de science ou environnement instrumentalisé : elle donne accès à d'autres dimensions de l'être. & « Il ne s'agit pas de\ldots{} mais tout simplement » ; chute \\
\end{longtable}

\textbf{Le thème et les œuvres.} Le thème de l'année est l'expérience de la nature. Le rapport 2026 rappelle que, pour la filière ATS, « seules deux œuvres sont étudiées » : Jules Verne, \emph{Vingt mille lieues sous les mers}, et Georges Canguilhem, \emph{La Connaissance de la vie}, dont le programme retient l'introduction, la partie I et les chapitres 2 à 5 de la partie III. Ce corrigé s'appuie sur l'introduction « La pensée et le vivant », sur « L'expérimentation en biologie animale » (partie I) et, dans la partie « Philosophie », sur « Machine et organisme », « Le vivant et son milieu » et « Le normal et le pathologique ». Le jury juge « de loin préférable » de s'en tenir à ces deux œuvres ; d'autres références restent possibles, mais en appoint.

\textbf{Difficulté.} Le résumé est classique mais dense : le texte est abstrait et ses idées se recoupent, ce qui pousse à recopier. La dissertation est abordable, car la citation est extraite du texte que l'on vient de résumer ; le piège est de s'en tenir à un « oui » enthousiaste ou à un « non » artificiel, alors que la phrase contient elle-même ses réserves (« bénéfique », « dans certains cas », « véritable »).

\textbf{Ce qui a départagé les copies.} Le rapport 2026 ne donne ni moyenne ni taux de réussite. Il juge les dissertations très hétérogènes et désigne ce qui a fait la différence : la connaissance précise des deux œuvres, l'analyse fine de la citation, le refus de plaquer le cours, l'équilibre des parties. Pour le résumé, il signale que le nombre de mots a été respecté cette année ; l'écart s'est donc fait sur la fidélité au plan, l'exhaustivité, la reformulation et la correction de la langue.

\textbf{À ne pas laisser passer.} Les points de contrat du résumé (fourchette, barres, total), qui se vérifient en quelques secondes ; la réponse explicite à « Partagez-vous ce point de vue ? » ; des références localisables (un épisode, un chapitre, un personnage) dans chaque sous-partie.

\textbf{Gestion du temps.} Le rapport 2026 conseille de commencer par le résumé, d'y passer au plus 1 h 30, et de garder 2 h 30 pour la dissertation, afin d'écrire au moins quatre pages en interligne double, puis de se réserver une relecture. Répartition possible :

\begin{itemize}
\tightlist
\item
  lecture et découpage du texte, 20 min ;
\item
  résumé (brouillon, réduction, recopie, décompte), 1 h 05 ;
\item
  analyse du sujet, problématique et plan détaillé avec les références, 40 min ;
\item
  rédaction, 1 h 35 (introduction et conclusion préparées au brouillon) ;
\item
  relecture, 20 min.
\end{itemize}

\hypertarget{partie-1-ruxe9sumuxe9}{%
\section{Partie 1 --- Résumé}\label{partie-1-ruxe9sumuxe9}}

Il s'agit de restituer le mouvement d'un texte argumentatif en le contractant d'environ un cinquième, sans commentaire, dans le système d'énonciation de l'auteur, en respectant un contrat formel strict.

\hypertarget{question-1-uxe9tape-a-comprendre-et-duxe9couper-le-texte}{%
\subsection{Question 1, étape a --- Comprendre et découper le texte}\label{question-1-uxe9tape-a-comprendre-et-duxe9couper-le-texte}}

\begin{indication}{1}

Repérez les mots de liaison qui ouvrent chaque paragraphe et ceux qui organisent le deuxième : ils dessinent le trajet de l'argumentation.

\end{indication}

\begin{indication}{2}

Demandez-vous quelle phrase porte la thèse et quelle phrase la reprend à la fin : l'une est au cœur du § 2, l'autre clôt le § 4. Ce qui les sépare (le § 3) est un élargissement historique.

\end{indication}

\hypertarget{solution}{%
\subsubsection{Solution}\label{solution}}

Le texte part d'un constat collectif (l'écologie engage chacun), le renverse (« Pourtant ») en montrant que le rapport vécu à la Nature s'est raréfié et qu'on sait mal ce qu'il produit chez l'individu, puis énonce sa thèse : l'expérience privilégiée de la Nature touche l'être entier, change le regard, \textbf{surtout} reconduit à soi et parfois transforme. Le § 3 replace cette thèse dans l'histoire des représentations occidentales (amour et exploitation, maîtrise contre appartenance, recul de la maîtrise). Le § 4 en tire une figure, l'habitant de la Terre qui comprend les interdépendances dont il fait partie, et une chute qui reprend la thèse : la Nature ouvre l'homme à d'autres dimensions de son être. Le tableau de la présentation donne le détail du découpage.

Proportions à respecter : les § 2 et 4, qui portent la thèse, méritent plus de la moitié du résumé ; les exemples (tourisme vert, sports de nature) et la référence à Lussault disparaissent en tant que tels, seule leur idée subsiste.

\begin{piege}

Consacrer l'essentiel du budget aux deux premiers paragraphes, les plus concrets, et expédier la fin en une phrase. Un résumé qui ne dit pas que la Nature \emph{ouvre l'homme à d'autres dimensions de son être} a perdu la conclusion du texte.

\end{piege}

\begin{jury}

Le rapport 2026 rappelle que « la composition du résumé doit tenir compte du plan du texte initial » et que la restitution doit être exhaustive : « Les résumés lacunaires sont fortement pénalisés. » Il demande aussi de rester fidèle à l'extrait, sans idée extérieure, et d'écrire un résumé autonome, compréhensible sans le texte.

\end{jury}

\hypertarget{question-1-uxe9tape-b-fixer-la-convention-et-le-budget}{%
\subsection{Question 1, étape b --- Fixer la convention et le budget}\label{question-1-uxe9tape-b-fixer-la-convention-et-le-budget}}

\begin{indication}{1}

La fourchette est de 108 à 132 mots. Avant d'écrire, décidez comment vous comptez « l'homme », « qu'il » ou « lui-même », et visez une valeur qui reste dans la fourchette même si le correcteur compte autrement.

\end{indication}

\begin{indication}{2}

Répartissez environ 120 mots entre les quatre mouvements en fonction de leur poids dans l'argumentation, pas de leur longueur dans le texte.

\end{indication}

\hypertarget{solution-1}{%
\subsubsection{Solution}\label{solution-1}}

Convention retenue, la plus répandue dans les concours : est un mot toute unité séparée par un blanc ou par une apostrophe d'élision (« l'homme » = 2 mots, « qu'il » = 2 mots) ; « aujourd'hui » et les mots composés à trait d'union (« lui-même ») comptent pour un mot. On l'indique brièvement avec le total.

Viser environ 125 mots selon cette convention garantit d'être aussi dans la fourchette pour un correcteur qui compterait entre blancs (environ 115).

Budget indicatif : § 1, une douzaine de mots ; § 2, une cinquantaine (constat de rareté, invention individuelle des natures, thèse en trois temps) ; § 3, une trentaine ; § 4, une trentaine (habitant de la Terre, chute).

\begin{methode}

Fixez la convention au brouillon et comptez sur la version recopiée, par tranches de vingt : chaque barre posée est un point de contrôle, et une erreur se voit tout de suite.

\end{methode}

\begin{jury}

Rapport 2026 : « il est souhaitable de s'approcher de la marge maximale autorisée, soit 132 mots, mais tout dépassement est sanctionné, surtout si le jury constate une fraude sur le décompte du nombre de mots ». Viser environ 125 mots selon la convention usuelle concilie les deux exigences : on reste près du plafond sans risquer de le franchir si le correcteur compte autrement.

\end{jury}

\hypertarget{question-1-uxe9tape-c-ruxe9diger-compter-poser-les-barres}{%
\subsection{Question 1, étape c --- Rédiger, compter, poser les barres}\label{question-1-uxe9tape-c-ruxe9diger-compter-poser-les-barres}}

\begin{indication}{1}

Rédigez une phrase par unité de sens, au présent et sans « l'auteur dit que », puis élaguez : adjectifs redondants, exemples, doublets.

\end{indication}

\begin{indication}{2}

Conservez les connecteurs qui portent la logique : l'opposition du § 2, la gradation de la thèse, le refus suivi d'une proposition au § 4, la conclusion. Vérifiez que chaque connecteur rend bien la relation du texte.

\end{indication}

\hypertarget{solution-2}{%
\subsubsection{Solution}\label{solution-2}}

\textbf{Problématique du texte :} au-delà de l'enjeu collectif de l'écologie, que fait à l'individu l'expérience vécue de la Nature ?

Les enjeux écologiques obligent chacun à interroger sa relation à la Nature. Or nos sociétés technicisées ne la côtoient plus \textbar{} que rarement, indirectement : on tente de renouer, mais le vécu individuel échappe. Chacun, comme la société, invente ses natures, éprouvées \textbar{} ailleurs grâce aux voyages. Engageant l'être entier, cette expérience renouvelle le regard, mais surtout reconduit à soi, voire transforme. \textbar{}

L'Occident aime et détruit la Nature ; l'idée d'un homme souverain, ruinée par ses ravages et le brouillage \textbar{} de la frontière animale, cède devant celle d'une espèce solidaire du vivant.

Mieux qu'un militant, il faudrait un \textbar{} habitant conscient des interdépendances qui l'unissent au monde. Car, irréductible au simulacre, au savoir, au politique, la Nature ouvre \textbar{} l'homme à d'autres dimensions de lui-même.

\textbf{Total : 128 mots} (décompte usuel : le mot élidé compte pour un mot, « l'homme » = 2 mots ; un mot composé à trait d'union compte pour un mot).

Vérification par tranches (convention usuelle) :

\begin{longtable}[]{@{}lll@{}}
\toprule\noalign{}
Tranche & Fin de la tranche & Mots \\
\midrule\noalign{}
\endhead
\bottomrule\noalign{}
\endlastfoot
1 & \ldots{} ne la côtoient plus & 20 \\
2 & \ldots{} invente ses natures, éprouvées & 20 \\
3 & \ldots{} mais surtout reconduit à soi, voire transforme. & 20 \\
4 & \ldots{} ruinée par ses ravages et le brouillage & 20 \\
5 & \ldots{} Mieux qu'un militant, il faudrait un & 20 \\
6 & \ldots{} au savoir, au politique, la Nature ouvre & 20 \\
7 & l'homme à d'autres dimensions de lui-même. & 8 \\
Total & & 128 \\
\end{longtable}

Compté entre blancs, sans séparer les élisions, le même résumé fait 119 mots : il reste dans la fourchette.

Les trois paragraphes suivent le plan du texte : les § 1 et 2 (constat, puis thèse), le § 3 (le paradoxe occidental), le § 4 (l'habitant de la Terre et la chute).

Correspondances principales : « côtoient plus que rarement, indirectement » rend « épisodiquement et d'une manière indirecte » ; « Chacun, comme la société, invente ses natures » rend « cela vaut aussi pour l'individu » ; « renouvelle le regard, mais surtout reconduit à soi, voire transforme » garde les trois temps de la thèse et sa gradation, « voire » traduisant « dans certains cas » ; « Mieux qu'un militant » rend le refus de l'« écocitoyen » ; « irréductible au simulacre, au savoir, au politique » condense la triple négation finale.

\begin{piege}

Fausser un connecteur en voulant gagner des mots. Écrire « Chacun invente pourtant ses natures » introduit une opposition que le texte ne fait pas : il dit que l'invention des natures vaut \emph{aussi} pour l'individu.

\end{piege}

\begin{piege}

Recopier des groupes de mots du texte (« du sensoriel au transcendant », « habitant de la Terre », « tourisme vert et sports de nature ») : le résumé doit reformuler, et les exemples n'y ont pas leur place. Le rapport 2026 admet seulement qu'on reprenne les concepts fondamentaux qui ne se reformulent pas, comme « Nature » ici.

\end{piege}

\begin{piege}

Gagner des mots en supprimant le verbe principal. Le rapport 2026 est net : « Un groupe nominal articulé autour d'un participe présent ne saurait constituer une phrase recevable. » « Engageant l'être entier, cette expérience renouvelle le regard » est une phrase ; « Une expérience engageant l'être entier » n'en est pas une.

\end{piege}

\begin{jury}

Sur la présentation, le rapport 2026 tolère « à la limite le paragraphe monobloc », mais juge inadmissibles les « productions constituées de sept ou huit phrases avec retour à la ligne ». Quelques paragraphes calqués sur les mouvements du texte rendent la structure visible. Il rappelle aussi qu'un résumé truffé de fautes de syntaxe « ne peut prétendre à une note supérieure à 1/10 ».

\end{jury}

\hypertarget{partie-2-dissertation}{%
\section{Partie 2 --- Dissertation}\label{partie-2-dissertation}}

La citation est extraite du texte résumé ; il faut la discuter à l'aide des œuvres, en montrant à quelles conditions l'expérience de la Nature transforme réellement celui qui l'éprouve.

\hypertarget{question-2-uxe9tape-a-analyser-la-citation}{%
\subsection{Question 2, étape a --- Analyser la citation}\label{question-2-uxe9tape-a-analyser-la-citation}}

\begin{indication}{1}

Prenez la phrase mot à mot : quels termes posent une condition, lesquels hiérarchisent, lesquels portent un jugement de valeur ?

\end{indication}

\begin{indication}{2}

Interrogez particulièrement « privilégié », « initie », les guillemets de « façon », la construction « non seulement\ldots{} mais surtout », puis « bénéfique », « dans certains cas » et « véritable ».

\end{indication}

\hypertarget{solution-3}{%
\subsubsection{Solution}\label{solution-3}}

\begin{longtable}[]{@{}
  >{\raggedright\arraybackslash}p{(\columnwidth - 4\tabcolsep) * \real{0.2348}}
  >{\raggedright\arraybackslash}p{(\columnwidth - 4\tabcolsep) * \real{0.3493}}
  >{\raggedright\arraybackslash}p{(\columnwidth - 4\tabcolsep) * \real{0.3860}}@{}}
\toprule\noalign{}
\begin{minipage}[b]{\linewidth}\raggedright
Terme
\end{minipage} & \begin{minipage}[b]{\linewidth}\raggedright
Sens
\end{minipage} & \begin{minipage}[b]{\linewidth}\raggedright
Prise pour la discussion
\end{minipage} \\
\midrule\noalign{}
\endhead
\bottomrule\noalign{}
\endlastfoot
« rapport privilégié » & relation rare, intense, réservée & un privilège suppose des conditions (moyens, technique, circonstances) : le rapport n'est pas donné à tous \\
« initie » & introduire, faire franchir un seuil & vocabulaire de l'initiation : un seuil, un guide, des épreuves, un avant et un après \\
« une autre ``façon'' de la voir » & une autre manière de voir ; les guillemets suggèrent que c'est la forme même donnée à l'objet qui change & voir autrement n'est pas voir plus, c'est former autrement ce que l'on voit (lecture interprétative, à présenter avec prudence) \\
« non seulement\ldots{} mais surtout » & gradation & le gain de connaissance est subordonné au gain pour le sujet : c'est le point le plus discutable \\
« rencontre avec soi-même », « retour sur soi » & réflexivité & rencontre-t-on soi-même, ou sa condition de vivant ? \\
« bénéfique » & jugement de valeur & rien ne garantit que le retour sur soi soit bon : il peut enfermer \\
« dans certains cas » & restriction & la transformation est possible, non nécessaire : à quelles conditions ? \\
« véritable transformation » & changement de ce que l'on est, non d'humeur & il faut un critère pour distinguer une transformation véritable d'une émotion forte \\
\end{longtable}

\begin{methode}

Les restrictions qu'une citation contient elle-même (« dans certains cas », « véritable ») sont les meilleures prises pour la discussion : elles permettent de nuancer sans contredire artificiellement.

\end{methode}

\begin{jury}

Le rapport 2026 demande, dans l'introduction, de « citer intégralement le sujet, puis analyser les notions et concepts importants ». Il a valorisé les copies « qui ont repéré la granularité du sujet », c'est-à-dire qui ont pesé chacun de ses termes.

\end{jury}

\hypertarget{question-2-uxe9tape-b-probluxe9matiser}{%
\subsection{Question 2, étape b --- Problématiser}\label{question-2-uxe9tape-b-probluxe9matiser}}

\begin{indication}{1}

Cherchez la tension interne de la thèse : qu'est-ce qui fait la valeur d'une expérience de la Nature, et qu'advient-il de cette valeur si l'essentiel est le retour sur soi ?

\end{indication}

\begin{indication}{2}

Relisez la fin du texte support : la Nature y est dite « pas uniquement » ersatz, objet de science ou environnement instrumentalisé. Que devient-elle si elle sert surtout à se découvrir ?

\end{indication}

\hypertarget{solution-4}{%
\subsubsection{Solution}\label{solution-4}}

Ce qui fait le prix d'une expérience de la Nature, c'est son altérité : elle nous confronte à ce qui n'est pas nous. Si son bénéfice essentiel est un retour sur soi, la Nature risque de n'être plus qu'un miroir, un moyen, ce que le texte lui-même refuse de réduire à ses usages.

\textbf{Problématique :} le rapport privilégié à la Nature nous reconduit-il à nous-mêmes, ou ne nous transforme-t-il qu'en nous délogeant de la place que nous croyions occuper ?

\begin{piege}

Transformer la question en débat écologique (réchauffement, biodiversité, engagement). Le sujet porte sur ce que l'expérience de la Nature fait à celui qui l'éprouve ; le texte écarte d'ailleurs explicitement la figure de l'« écocitoyen ».

\end{piege}

\begin{jury}

Selon le rapport 2026, le sujet « n'était guère traitable sans avoir au préalable travaillé sur le texte », qui donnait des clés pour comprendre la citation et des arguments pour la discuter. Le résumé fait d'abord sert donc la dissertation : la tension entre altérité et retour sur soi vient du dernier paragraphe du texte.

\end{jury}

\hypertarget{question-2-uxe9tape-c-construire-le-plan-et-choisir-les-ruxe9fuxe9rences}{%
\subsection{Question 2, étape c --- Construire le plan et choisir les références}\label{question-2-uxe9tape-c-construire-le-plan-et-choisir-les-ruxe9fuxe9rences}}

\begin{indication}{1}

Un plan en trois temps convient : accorder ce que la thèse a de juste, montrer ses limites à partir de ses propres restrictions, puis requalifier le mot le plus vague de la citation.

\end{indication}

\begin{indication}{2}

Chez Verne, comparez ce que le même voyage fait aux quatre hommes du Nautilus. Chez Canguilhem, cherchez un chapitre qui pense le rapport du vivant à son milieu, et un autre qui fournit un critère pour distinguer un vrai changement d'un simple état passager.

\end{indication}

\hypertarget{solution-5}{%
\subsubsection{Solution}\label{solution-5}}

\begin{longtable}[]{@{}
  >{\raggedright\arraybackslash}p{(\columnwidth - 6\tabcolsep) * \real{0.2018}}
  >{\raggedright\arraybackslash}p{(\columnwidth - 6\tabcolsep) * \real{0.2245}}
  >{\raggedright\arraybackslash}p{(\columnwidth - 6\tabcolsep) * \real{0.2782}}
  >{\raggedright\arraybackslash}p{(\columnwidth - 6\tabcolsep) * \real{0.2656}}@{}}
\toprule\noalign{}
\begin{minipage}[b]{\linewidth}\raggedright
Partie
\end{minipage} & \begin{minipage}[b]{\linewidth}\raggedright
Idée directrice
\end{minipage} & \begin{minipage}[b]{\linewidth}\raggedright
Verne
\end{minipage} & \begin{minipage}[b]{\linewidth}\raggedright
Canguilhem
\end{minipage} \\
\midrule\noalign{}
\endhead
\bottomrule\noalign{}
\endlastfoot
I. Oui : une initiation qui fait retour sur le sujet & 1. rupture du regard ; 2. structure d'initiation ; 3. retour sur soi & salon et panneaux ; forêt de l'île Crespo ; Conseil classificateur et Ned Land ; seuil, Nemo guide, « La mer est tout », Atlantide, banquise, poulpes, \emph{Mobilis in mobili} ; Aronnax hésite à fuir & introduction : la connaissance décompose ce que la vie compose ; « tenir du vivant l'idée du vivant » ; « L'expérimentation en biologie animale » : obstacles propres au vivant \\
II. Mais : ni nécessaire, ni toujours bénéfique, risque du miroir & 1. Conseil et Ned Land ; 2. Nemo enfermé dans sa douleur ; 3. la Nature réduite à un miroir & les trois passagers ; rupture avec la société, cimetière de corail, pêcheur de perles, éperonnage, portrait, « Dieu tout-puissant ! assez ! assez ! » ; pavillon au pôle Sud puis banquise & « Le vivant et son milieu » : \emph{Umgebung}, \emph{Umwelt}, \emph{Welt}, la tique \\
III. En vérité : rencontrer sa condition de vivant ; décentration médiatisée ; critère de normativité & 1. se découvrir vivant parmi les vivants ; 2. médiation de la technique et du récit ; 3. transformation = gain de normativité & monde qui existait sans témoin ; le Nautilus ; le récit d'Aronnax ; sort comparé des personnages ; réponse à l'Ecclésiaste & « Le vivant et son milieu » (tout vivant est centre) ; « Machine et organisme » ; « Le normal et le pathologique » \\
\end{longtable}

\begin{methode}

Dressez ce tableau au brouillon avant de rédiger : il garantit que chaque œuvre apparaît dans chaque partie et qu'un même épisode ne sert pas deux fois.

\end{methode}

\begin{jury}

Le rapport 2026 n'impose pas de plan type ; deux parties suffisent, et de bons plans en deux parties ont eu la note maximale. Il exige en revanche que la pensée reste cohérente : « les éléments d'antithèse ne peuvent venir que nuancer ce qui a pu être dit précédemment ». Ici, la partie II limite la thèse sans la nier. Le rapport ajoute : « On évitera enfin de limiter l'étude d'une œuvre à une partie », car les points communs et les divergences des auteurs doivent être examinés.

\end{jury}

\begin{piege}

Déséquilibrer le devoir. Le rapport 2026 déplore « une première partie pléthorique » suivie de sections « croupions », et sanctionne les copies qui « plaquent des paragraphes entiers du cours » sur les auteurs au lieu de traiter la citation. Chaque référence du tableau doit servir une idée de la partie où elle figure.

\end{piege}

\begin{piege}

Attribuer à Verne les enthousiasmes ou les jugements de ses personnages. Aronnax est le narrateur, Nemo un personnage : confronter ce qu'ils disent à ce qu'ils font fournit les meilleurs arguments.

\end{piege}

\begin{piege}

Affirmer des faits incertains sur une œuvre. Le roman ne dit pas que Nemo meurt dans le Maelström : Aronnax ignore si le Nautilus a résisté. Une erreur sur le dénouement ruine l'argument qui s'y appuie.

\end{piege}

\hypertarget{question-2-uxe9tape-d-ruxe9diger-lintroduction-et-la-conclusion}{%
\subsection{Question 2, étape d --- Rédiger l'introduction et la conclusion}\label{question-2-uxe9tape-d-ruxe9diger-lintroduction-et-la-conclusion}}

\begin{indication}{1}

L'introduction enchaîne une amorce précise (de préférence tirée d'une œuvre), la citation recopiée exactement, son analyse, la tension, la problématique et l'annonce du plan.

\end{indication}

\begin{indication}{2}

La conclusion répond à la question posée (« oui, à condition que\ldots{} ») en reprenant le trajet, puis ouvre sur une question issue du devoir lui-même.

\end{indication}

\hypertarget{solution-6}{%
\subsubsection{Solution}\label{solution-6}}

Voir l'introduction et la conclusion du devoir complet à l'étape suivante. L'amorce part du salon du Nautilus, ce qui place d'emblée une œuvre au programme au cœur du sujet. Avant l'annonce du plan, l'introduction nomme les deux œuvres, avec leurs auteurs et leurs dates. La conclusion répond explicitement (« Nous partageons donc ce point de vue, à condition d'en déplacer le centre ») et ouvre sur la transmission de l'expérience par le récit.

\begin{jury}

Le rapport 2026 détaille l'introduction attendue : une accroche sans banalité, puis la citation et son analyse, et il faut « rappeler que le sujet sera traité à la lumière des deux œuvres au programme (qu'il convient de citer explicitement) », avant le problème et l'annonce du plan. La conclusion « est indispensable » : elle répond « de manière claire à la problématique posée en introduction », rappelle brièvement le parcours et ouvre sans lieu commun.

\end{jury}

\begin{piege}

Oublier de nommer les œuvres dans l'introduction parce que l'accroche en cite déjà une. L'accroche n'en dispense pas : le rapport 2026 demande de citer explicitement les deux œuvres au moment de poser le problème.

\end{piege}

\hypertarget{question-2-uxe9tape-e-ruxe9diger-le-devoir-complet}{%
\subsection{Question 2, étape e --- Rédiger le devoir complet}\label{question-2-uxe9tape-e-ruxe9diger-le-devoir-complet}}

\begin{indication}{1}

Chaque sous-partie s'ouvre sur une idée, la démontre à l'aide d'au moins une référence précise, et se termine par une phrase qui la rattache au sujet.

\end{indication}

\begin{indication}{2}

Faites dialoguer les œuvres : une formule de Canguilhem doit expliquer un épisode de Verne, et non s'y ajouter comme une illustration de plus.

\end{indication}

\hypertarget{solution-7}{%
\subsubsection{Solution}\label{solution-7}}

\textbf{Problématique :} le rapport privilégié à la Nature nous reconduit-il à nous-mêmes, ou ne nous transforme-t-il qu'en nous délogeant de la place que nous croyions occuper ?

\textbf{Introduction}

Lorsque, dans \emph{Vingt mille lieues sous les mers}, les panneaux du salon du Nautilus s'écartent pour la première fois sur la mer illuminée par le fanal du navire, le professeur Aronnax, qui avait consacré un ouvrage savant aux grands fonds sous-marins, mesure soudain la distance entre ce qu'il savait et ce qu'il voit. Le spectacle ne lui apprend pas seulement des espèces : il ébranle le savant lui-même. C'est une expérience de ce genre que décrivent Stéphanie Chanvallon et Stéphane Héas lorsqu'ils affirment qu'« un rapport privilégié avec la Nature non seulement initie à une autre ``façon'' de la voir, mais surtout constitue une rencontre avec soi-même, un retour sur soi bénéfique », et qu'« il permet dans certains cas une véritable transformation de soi ».

La phrase est construite avec soin. Le rapport dont il est question est dit \emph{privilégié} : ce n'est pas le commerce ordinaire avec un paysage, mais une relation rare et intense, qui suppose des conditions que tous ne réunissent pas. Le verbe \emph{initier} appartient au vocabulaire du seuil : on n'apprend pas la Nature comme une leçon, on y est introduit. Les guillemets qui entourent « façon » signalent que l'enjeu dépasse une simple manière de regarder : c'est la forme même que nous donnons à ce que nous voyons qui se trouve modifiée. Surtout, la gradation « non seulement\ldots{} mais surtout » établit une hiérarchie : le gain de connaissance est second, la rencontre avec soi est l'essentiel. Enfin, deux réserves encadrent la thèse : le retour sur soi est déclaré \emph{bénéfique}, ce qui est un jugement de valeur, et la transformation n'est promise que « dans certains cas ».

Or cette thèse recèle une difficulté. Ce qui fait le prix d'une expérience de la Nature, c'est qu'elle nous met en présence de ce qui n'est pas nous : une altérité, une résistance, une indifférence. Si l'essentiel de la rencontre est un retour sur soi, la Nature risque de n'être plus qu'un miroir, un moyen de se découvrir, alors même que les auteurs rappellent, à la fin de leur texte, qu'elle « n'est pas uniquement » une nature recréée, un objet de science ou un environnement instrumentalisé.

Le rapport privilégié à la Nature nous reconduit-il donc à nous-mêmes, ou ne nous transforme-t-il qu'en nous délogeant de la place que nous croyions occuper ? Pour en juger, on confrontera les deux œuvres au programme : \emph{Vingt mille lieues sous les mers} (1869-1870), où Jules Verne fait vivre à quatre hommes une même immersion dans l'océan, et \emph{La Connaissance de la vie} (1952), où Georges Canguilhem pense le rapport du vivant à son milieu et celui du savant au vivant qu'il étudie. Nous verrons d'abord que l'expérience privilégiée de la Nature a bien la structure d'une initiation, qui change le regard et fait retour sur celui qui l'éprouve. Nous montrerons ensuite que ce retour sur soi n'a rien de nécessaire, qu'il n'est pas toujours bénéfique et qu'il peut réduire la Nature au rang de miroir. Nous soutiendrons enfin que ce que l'on rencontre est moins un « soi » que sa condition de vivant, et que la transformation n'est véritable que si cette décentration ouvre à de nouvelles manières de vivre.

\textbf{I. Le rapport privilégié à la Nature est bien une initiation qui change le regard et reconduit le sujet à lui-même}

\textbf{1. Une rupture du regard.} Il faut d'abord accorder aux auteurs que le rapport privilégié commence par une rupture avec le regard ordinaire, qui fait de la Nature un décor ou une ressource. Chez Verne, cette rupture a un lieu précis : le salon du Nautilus, dont les panneaux ouverts transforment la mer en un spectacle que nul n'avait jamais contemplé ainsi. La première promenade sous-marine, dans la forêt de l'île Crespo, l'accomplit : revêtus de scaphandres, Aronnax et Conseil marchent au fond de l'eau, dans une lumière et un silence inconnus, et la Nature cesse d'être un objet placé devant eux pour devenir un milieu dans lequel ils se tiennent. Le roman souligne d'ailleurs, sur le mode comique, l'insuffisance des cadres habituels : Conseil, classificateur infatigable, range chaque poisson dans son embranchement, sa classe et son ordre, tandis que Ned Land, qui les connaît en pêcheur, les nomme sans savoir les classer ; aucun des deux savoirs n'épuise le spectacle.

Canguilhem permet de comprendre ce que signifie voir « autrement ». Dans l'introduction de \emph{La Connaissance de la vie}, « La pensée et le vivant », il montre que la connaissance procède en décomposant ce que la vie compose : elle analyse, sépare, immobilise, alors que le vivant forme des totalités. Voir autrement la Nature, ce n'est donc pas obtenir une image plus fine du même objet, c'est suspendre, au moins un instant, l'attitude qui découpe pour accueillir ce qui se présente comme un tout. C'est en ce sens que la « façon » de voir est aussi une façon de \emph{former} son objet.

\textbf{2. La structure d'une initiation.} Ce déplacement mérite le nom d'initiation, car il en a les trois éléments : un seuil, un maître, des épreuves. Le seuil, chez Verne, est franchi malgré soi : projeté à la mer par le choc qui ébranle l'\emph{Abraham Lincoln}, Aronnax se retrouve avec Conseil et Ned Land enfermé dans une cellule obscure, puis soumis à la loi du capitaine, qui leur interdit tout retour à terre. Le maître est Nemo, qui choisit ce que ses hôtes verront, leur ouvre ou leur ferme les panneaux, et formule le credo de cette vie nouvelle : « La mer est tout ». Les épreuves enfin rythment le voyage : la marche nocturne parmi les ruines de l'Atlantide, à la lueur d'un volcan sous-marin ; l'emprisonnement sous la banquise ; le combat contre les poulpes géants. Chaque épisode est un seuil franchi. La devise du Nautilus, \emph{Mobilis in mobili}, « mobile dans l'élément mobile », résume cette logique : il ne s'agit plus de juger la mer depuis la terre ferme, mais d'épouser son mouvement.

\textbf{3. Un retour sur soi.} Reste le point décisif de la citation : ce déplacement fait retour sur le sujet. Canguilhem en donne le fondement. Celui qui connaît la vie n'est jamais extérieur à ce qu'il connaît : le biologiste est lui-même un vivant, et sa connaissance est une manière pour la vie de se rapporter à elle-même. D'où la formule par laquelle s'achève son introduction : la pensée du vivant doit « tenir du vivant l'idée du vivant ». Le chapitre « L'expérimentation en biologie animale » le confirme concrètement : la spécificité des espèces, l'individualité des organismes, leur caractère de totalité, l'irréversibilité de leurs processus font obstacle à l'expérimentation, et le savant découvre, en se heurtant au vivant, les limites de son propre pouvoir de connaître. Connaître le vivant, c'est donc aussi apprendre quelque chose sur soi.

Verne montre ce retour sur soi dans la trajectoire d'Aronnax. Le savant observe, décrit, classe, mais il éprouve aussi un enthousiasme qui lui fait oublier sa captivité, au point qu'il hésite à suivre Ned Land dans ses projets d'évasion, tant le voyage le retient. Sous le professeur du Muséum apparaît un homme que la mer captive. À ce premier niveau, on peut donc partager le point de vue des auteurs : le rapport privilégié à la Nature initie, change la façon de voir et révèle à lui-même celui qui voit.

\textbf{II. Mais ce retour sur soi n'est ni nécessaire ni toujours bénéfique, et il risque de faire de la Nature un miroir}

\textbf{1. Une transformation qui n'a rien d'automatique.} Les auteurs le reconnaissent eux-mêmes : la transformation n'a lieu que « dans certains cas ». Verne l'illustre en faisant vivre le même voyage à trois passagers. Conseil, serviteur dévoué, traverse les mêmes merveilles qu'Aronnax sans cesser d'être celui qu'il était : il classe, il suit son maître, et sa science range sans rencontrer. Ned Land, le harponneur canadien, vit le voyage comme une captivité : il réclame de la viande fraîche, guette chaque terre en vue et n'aspire qu'à s'enfuir. Le même privilège, que nul autre homme ne partage, produit trois effets incomparables. Il n'existe donc pas de rapport à la Nature qui agirait uniformément sur tous ; l'expérience dépend de ce que chacun y cherche et de ce qu'il consent à y laisser.

\textbf{2. Un retour sur soi qui peut enfermer.} Plus grave, le retour sur soi n'est pas nécessairement bénéfique. Nemo possède le rapport le plus privilégié qui se puisse concevoir avec la mer. Il a rompu avec les hommes --- « Je ne suis pas ce que vous appelez un homme civilisé ! J'ai rompu avec la société tout entière pour des raisons que moi seul j'ai le droit d'apprécier » --- et il tire de l'océan sa nourriture, ses vêtements, son énergie, jusqu'à la sépulture de ses compagnons, inhumés dans un cimetière de corail. Or cette vie rendue à la Nature ne l'apaise pas. Sa compassion même, lorsqu'il sauve un pêcheur de perles de Ceylan, le ramène aux peuples opprimés dont il se dit solidaire, c'est-à-dire à sa blessure. Et sa haine reste intacte : il éperonne un navire de guerre et le regarde sombrer, avant de s'agenouiller en sanglotant devant le portrait d'une femme et de deux enfants. Les dernières paroles qu'Aronnax entend de lui, au moment de la fuite, sont un cri : « Dieu tout-puissant ! assez ! assez ! ». La mer n'a pas transformé Nemo ; elle lui a donné les moyens de rester enfermé dans sa douleur. Se rencontrer soi-même, c'est aussi risquer de s'y installer, et la solitude au sein de la Nature devient alors un retranchement.

\textbf{3. Le risque de faire de la Nature un miroir.} La dernière réserve vise la structure de la thèse. Si l'essentiel du rapport à la Nature est le bénéfice que le sujet en retire, la Nature vaut d'abord comme moyen. Les auteurs écrivent pourtant qu'elle « n'est pas uniquement » une nature « que l'on recrée en ersatz », ni « un objet de science », ni « cet environnement instrumentalisé » : faire d'elle avant tout un révélateur de soi, ce serait lui assigner un usage de plus, d'autant plus discret qu'il se présente comme une forme de respect.

Canguilhem aide à formuler cette limite. Dans « Le vivant et son milieu », il reprend la distinction établie par le biologiste Uexküll entre l'\emph{Umgebung}, l'environnement géographique, l'\emph{Umwelt}, le milieu de comportement propre à un vivant, découpé dans cet environnement selon ce qui a pour lui une signification vitale (la tique ne retient du monde que quelques signaux), et la \emph{Welt}, l'univers de la science. Chaque vivant est le centre de son milieu propre. Une expérience de la Nature qui ne ferait que revenir au sujet resterait prisonnière de ce centre : elle ne verrait de la Nature que ce qui intéresse l'homme.

Verne met en scène cette limite avec ironie. Parvenu au pôle Sud, Nemo déploie un pavillon noir marqué d'un N d'or et prend possession de cette terre : le rapport le plus intime avec la Nature s'achève en geste de souveraineté, celui même de « l'Homme, être supérieur et maître de la Nature » que décrit le texte. Or peu après, le Nautilus est pris sous la banquise, l'air vient à manquer, et l'équipage doit attaquer la glace pour survivre. La mer n'est ni le miroir ni la complice de l'homme ; elle ignore les significations dont il l'investit. La citation n'est donc pas fausse, mais elle nomme mal ce qui est rencontré.

\textbf{III. Ce que l'on rencontre est moins soi-même que sa condition de vivant : la transformation véritable est une décentration}

\textbf{1. Se découvrir vivant parmi les vivants.} Il faut donc requalifier le « soi » de la citation. Ce que révèle l'expérience privilégiée de la Nature n'est pas d'abord une intériorité à retrouver, mais une condition : celle d'un vivant parmi d'autres. Canguilhem permet de le penser. Si chaque vivant est le centre de son milieu, aucun ne l'est absolument : la tique comme l'homme découpent leur monde selon leurs besoins. Découvrir que l'on est un centre parmi d'autres, c'est relativiser le sien. Le retour sur soi devient alors une décentration, et il n'est bénéfique que dans la mesure où il décentre.

C'est la leçon que reçoit Aronnax. Sous la mer, il découvre un monde qui existait sans lui : les récifs, les forêts sous-marines, les ruines de l'Atlantide englouties depuis des millénaires. Venu chasser un « monstre » à classer, il se découvre l'hôte d'un monde qui se passait de témoin. Le texte de Chanvallon et Héas ne dit pas autre chose dans son dernier mouvement, lorsqu'il souhaite un homme qui agirait « par compréhension des lois de la Nature », de ses « réseaux d'interdépendances », « de tout ce dont il fait partie ». La Nature ne renvoie pas l'homme à lui-même : elle l'inclut.

\textbf{2. Une décentration toujours médiatisée.} Cette décentration n'est jamais immédiate. Il n'y a pas d'expérience nue de la Nature, et c'est précisément ce qui rend le rapport « privilégié » : il suppose un dispositif. Chez Verne, ce dispositif est le Nautilus. Sans lui, ni forêt de Crespo, ni Atlantide, ni pôle Sud : l'océan resterait une surface. L'expérience la plus intime de la mer est donc l'œuvre d'une machine, mue par l'électricité. On pourrait y voir une contradiction, comme si la technique éloignait de la Nature ; Canguilhem invite à penser l'inverse. Dans « Machine et organisme », il refuse d'expliquer l'organisme par la machine et propose de comprendre la machine à partir de la vie : la technique prolonge l'activité du vivant au lieu de s'y opposer. Le Nautilus n'est pas l'écran qui masque la mer, mais l'organe par lequel un vivant terrestre accède à un milieu qui n'est pas le sien.

La médiation est aussi celle du récit. Le roman est le témoignage qu'Aronnax rédige après son retour : l'expérience ne devient compréhensible, et transmissible, qu'une fois mise en forme. Le texte support le suggérait en citant Lussault, pour qui « la société ``invente'' {[}\ldots{]} ses natures acceptables » : il n'y a pas de Nature en soi qui viendrait frapper une conscience vierge, mais des natures façonnées par des techniques, des savoirs et des récits.

\textbf{3. Un critère de la transformation véritable.} Reste à donner un sens précis à la « véritable transformation », faute de quoi n'importe quelle émotion forte passerait pour une métamorphose. Canguilhem en fournit le critère dans « Le normal et le pathologique » : la norme n'est pas une moyenne statistique, et le vivant est normatif, c'est-à-dire capable d'instituer des normes et d'en changer quand son milieu change. La santé n'est pas l'absence de troubles, mais la marge qui permet d'inventer d'autres manières de vivre. Une transformation est donc véritable lorsqu'elle accroît cette capacité : lorsqu'elle rend possible une manière de vivre et de penser qu'on ne concevait pas auparavant.

Ce critère éclaire les personnages de Verne. Conseil et Ned Land reviennent pour l'essentiel tels qu'ils étaient partis. Nemo s'est donné des normes propres, mais figées : sa haine se répète sans se convertir, et le roman l'abandonne dans le Maelström, sans que l'on sache s'il y a survécu. Aronnax seul revient autre. Les dernières lignes du roman le montrent répondant à la question de l'Ecclésiaste, « Qui a jamais pu sonder les profondeurs de l'abîme ? » : deux hommes ont désormais le droit d'y répondre, « le capitaine Nemo et moi ». Ce n'est pas une vanterie, mais la reconnaissance d'un savoir conquis par l'expérience, et qui demeure quand l'émotion est retombée. On peut alors répondre à la question posée : le rapport privilégié à la Nature transforme véritablement, non parce qu'il nous rendrait à nous-mêmes, mais parce qu'il nous ôte d'un centre que nous croyions occuper. C'est bien ce que suggèrent les auteurs lorsqu'ils préfèrent à l'« écocitoyen » « tout simplement un citoyen, un habitant de la Terre » : le mot juste n'est pas miroir, mais habitation.

\textbf{Conclusion}

La thèse de Chanvallon et Héas apparaît ainsi juste dans son constat et discutable dans sa formulation. Juste, car le rapport privilégié à la Nature a bien la structure d'une initiation : il rompt avec le regard ordinaire, engage celui qui l'éprouve, et Canguilhem rappelle qu'on ne connaît jamais la vie sans être soi-même vivant. Discutable, car le retour sur soi n'est ni automatique, comme le montrent Conseil et Ned Land, ni nécessairement bénéfique, comme le montre Nemo, et parce qu'en faisant de la rencontre avec soi l'essentiel, on risque de réduire la Nature au rôle de miroir. Nous partageons donc ce point de vue, à condition d'en déplacer le centre : ce qui transforme n'est pas le retour à soi mais la découverte de sa condition de vivant parmi les vivants, et cette transformation n'est véritable que si elle ouvre à de nouvelles normes de vie.

Une question demeure : si cette expérience est réservée à quelques privilégiés --- il n'y a qu'un Nautilus ---, comment pourrait-elle transformer d'autres qu'eux ? La réponse tient peut-être dans le geste d'Aronnax, qui raconte. Par le récit, l'expérience singulière devient commune, et \emph{Vingt mille lieues sous les mers} n'est pas seulement le récit d'une expérience de la nature : il en offre une à chacun de ses lecteurs.

\begin{methode}

Pour trouver une troisième partie, demandez-vous quel mot de la citation est le plus vague et requalifiez-le : ici, « soi-même » devient « condition de vivant ». Le dépassement reste ancré dans le sujet au lieu d'importer une idée étrangère.

\end{methode}

\begin{jury}

Le rapport 2026 définit le paragraphe comme une unité logique : une idée démontrée puis illustrée par les œuvres. Il ne faut pas passer à la ligne pour développer un exemple, mais « associer au sein d'une même unité graphique un argument et son illustration par un exemple ». Il relève aussi que les meilleures copies prennent en compte « la forme de l'œuvre, le geste d'écriture à l'origine des ouvrages au programme », notamment chez Verne : c'est le rôle du récit d'Aronnax en III.2.

\end{jury}

\begin{methode}

Les intertitres et les numéros en gras de ce corrigé servent la lecture à l'écran. Sur la copie, la structure se lit par les phrases d'annonce et de transition, et par la mise en page : le rapport 2026 recommande de ne sauter « une ligne ou deux qu'entre les parties ».

\end{methode}

\begin{piege}

Confondre normativité et normalité chez Canguilhem. Être normatif, c'est pouvoir instituer de nouvelles normes, non se conformer à une moyenne : écrire que l'expérience de la nature « rend plus normal » serait un contresens.

\end{piege}

\hypertarget{bilan}{%
\section{Bilan}\label{bilan}}

\begin{itemize}
\tightlist
\item
  \textbf{Résumé} : découper le texte en mouvements avant d'écrire ; fixer et indiquer une convention de décompte, et viser une longueur qui reste dans la fourchette sous les conventions courantes ; vérifier chaque connecteur ; ne conserver ni exemples ni syntagmes du texte ; compter sur la version recopiée, par tranches.
\item
  \textbf{Dissertation} : analyser les restrictions que la citation contient (« dans certains cas », « bénéfique », « véritable ») ; formuler la tension interne avant la problématique ; requalifier en troisième partie le terme le plus vague ; répondre explicitement à la question.
\item
  \textbf{Œuvres} : chez Verne, la comparaison des quatre hommes du Nautilus (Aronnax, Conseil, Ned Land, Nemo) répond à presque tous les sujets sur ce que l'expérience fait au sujet ; chez Canguilhem, le milieu propre (\emph{Umwelt}) et la normativité du vivant sont les deux outils les plus utiles. Revoir précisément le dénouement du roman et la lettre des quelques citations que l'on veut employer : le rapport 2026 conseille de relire les deux œuvres « en mémorisant quelques passages importants ».
\item
  \textbf{Ce qui départage} (rapport 2026) : une connaissance précise et personnelle des deux œuvres, une introduction qui analyse vraiment la citation et nomme les œuvres, des parties équilibrées qui font dialoguer les deux auteurs, un résumé fidèle au plan du texte et sans faute de syntaxe.
\item
  \textbf{Sujets transposables} : la nature comme miroir ou comme altérité ; expérience et transformation de soi ; rôle de la technique et du récit dans l'accès à la nature ; savoir et vécu.
\end{itemize}

\end{document}
